% Nouveau EGA V (cote 25): archive cote 25, site page 112 (PDF page 113).
% typescript. Draft compared: none.
% AI-transcribed by Claude Opus 5.5 (claude-opus-5-5, adaptive thinking, effort high), Message Batches API, batch ega25b (msgbatch_01A7n5YuXubUPNK1zsepi8RG), 2026-11-13; checked against both supplied scan parts; see CHECK_DECISIONS.tsv.
\scanpage{25/112}
\mspage{15}

\var{\deleted{\illegible{1}}}{}{deletion}{ega25_25_112-v0001} \var{\margin{\fbox{N$^{\text{os}}$ 3,4}}}{\placedmargin{\fbox{N$^{\text{os}}$ 3,4}}}{marginal note placed}{ega25_25_112-v0002} \hand{\unsure{3|2})} \var{\inserted{\uline{Systèmes linéaires de diviseurs, et morphismes dans les fibrés projectifs.}}}{\uline{Systèmes linéaires de diviseurs, et morphismes dans les fibrés projectifs.}}{insertion}{ega25_25_112-v0003}
\var{\deleted{16. \uline{Traductions en langage de systèmes linéaires.}}}{}{deletion}{ega25_25_112-v0004}

\var{\deleted{Soit X plat et loc de présentation finie sur S, et T un S-préschéma.}}{}{deletion}{ega25_25_112-v0005}
\var{\deleted{On appelle \uline{famille de} \var{\xout{\unsure{visieurs}}}{}{typewriter strike-out}{ega25_25_112-v0007} \uline{diviseurs} (sous-entendu: de Cartier)}}{}{deletion}{ega25_25_112-v0006}
\var{\deleted{\var{\margin{\unsure{r}}}{\placedmargin{\unsure{r}}}{marginal note placed}{ega25_25_112-v0009}\uline{positifs}\var{\xout{x}}{}{typewriter strike-out}{ega25_25_112-v0010}\hand{\unsure{$^{3}$|$^{z}$}} \var{\inserted[typed]{\uline{de X/S} \hand{,}}}{\uline{de X/S} \hand{,}}{insertion}{ega25_25_112-v0011} \hand{/} \uline{paramétrée par} S, un sous-préschéma fermé de $X\times_S T$ qui est}}{}{deletion}{ega25_25_112-v0008}
\var{\deleted{\var{\inserted{un diviseur de Cartier rel. à T, i.e. qui est}}{un diviseur de Cartier rel. à T, i.e. qui est}{insertion}{ega25_25_112-v0013} intersection complète \var{\xout{transversale}}{}{typewriter strike-out}{ega25_25_112-v0014} relative de codimension 1 relativement}}{}{deletion}{ega25_25_112-v0012}
\var{\deleted{à T en chacun de ses points, ou \var{\xout{\illegible{1}}}{}{typewriter strike-out}{ega25_25_112-v0016} ce qui revient au même, qui est plat}}{}{deletion}{ega25_25_112-v0015}
\var{\deleted{\var{\inserted[typed]{loc de prés finie}}{loc de prés finie}{insertion}{ega25_25_112-v0018} par rapport à T, et induit sur chaque fibre de $X_{T\,t} = X_{s\,k(t)}$ de $X_T$ sur}}{}{deletion}{ega25_25_112-v0017}
\var{\deleted{T un ``diviseur de Cartier positif'' i.e. un sous-préschéma qui est en}}{}{deletion}{ega25_25_112-v0019}
\var{\deleted{chaque \var{\inserted[typed]{de ses}}{de ses}{insertion}{ega25_25_112-v0021} point intersection complète \var{\inserted[typed]{dans $X_t$}}{dans $X_t$}{insertion}{ega25_25_112-v0022} de codimension 1. L'ensemble}}{}{deletion}{ega25_25_112-v0020}
\var{\deleted{$\mathrm{Div}_{X/S}(T)$ de telles familles de \var{\inserted[typed]{di}}{di}{insertion}{ega25_25_112-v0024}viseurs est, pour T variable, un}}{}{deletion}{ega25_25_112-v0023}
\var{\deleted{foncteur contravar\var{\xout{i}}{}{typewriter strike-out}{ega25_25_112-v0026}iant à valeurs dans (Ens), noté $\mathrm{Div}_{X/S}$ . Nous verrons}}{}{deletion}{ega25_25_112-v0025}
\var{\deleted{au Chap V qu'il est représentable lorsque $f:X\quad S$ est projectif. Lorsque}}{}{deletion}{ega25_25_112-v0027}
\var{\deleted{X = P(E) fibré projectif défini par E localement libre de type fini,}}{}{deletion}{ega25_25_112-v0028}
\var{\deleted{il est possible avec les seuls résultats du Chap III de prouver que}}{}{deletion}{ega25_25_112-v0029}
\var{\deleted{$\mathrm{Div}_{P/S}$ est représentable par le préschéma somme des $P(\mathrm{Symm}^n(\underline{E}))$ , les}}{}{deletion}{ega25_25_112-v0030}
\var{\deleted{composants de cette somme correspondant aux diviseurs de degré relatif}}{}{deletion}{ega25_25_112-v0031}
\var{\deleted{n , \var{\inserted{\unsure{cf|\S}\ n$^{\circ}$2}}{\unsure{cf|\S}\ n$^{\circ}$2}{insertion}{ega25_25_112-v0033} Ce résultat n'est pas indispensable pour la suite, et il suffira de}}{}{deletion}{ega25_25_112-v0032}
\var{\deleted{le signaler en remarque; nous l'inclurons peut-être quelque part dans}}{}{deletion}{ega25_25_112-v0034}
\var{\deleted{III, et au plus tard dans V .}}{}{deletion}{ega25_25_112-v0035}

Soit D une famille de diviseurs \var{\inserted[typed]{sur X/S}}{sur X/S}{insertion}{ega25_25_112-v0036} paramétrée par T (nous sous-enten%
% NOTE: the end of this line is cut off at the right edge of the scan
dons par la suite ``positif''). Un point $x\hand{\in} X$ est dit ``\uline{point fixe}'' \unsure{pour|pous}
cette famille de diviseurs si $pr_1^{-1}(x)\hand{\subset} D$ \var{\inserted{ensemblistement}}{ensemblistement}{insertion}{ega25_25_112-v0037}, de sorte que l'ensemble des
points non fixes est le complémentaire de $pr_1(\unsure{X}\times_S T - D)$; par suite, si
$T\hand{\longrightarrow} S$ est universellement ouvert (par exemple plat loc de prés finie)
l'ensemble des points fixes est fermé. On dit que la famille de diviseurs
est ``sans points fixes'' si l'ensemble \var{\inserted{$Z$}}{$Z$}{insertion}{ega25_25_112-v0038} des points fixes est vide.

\mspage{112}
